%%%%%%%%%%%%%%%%%%%%%%%%%%%%%%%%%%%%%%%%%%%%%%%%%%%%%%%%%%%%%%%%%%%%%%%%%%%%%%%%
% Obxectivo: Lista de termos empregados no documento,                          %
%            xunto cos seus respectivos significados.                          %
%%%%%%%%%%%%%%%%%%%%%%%%%%%%%%%%%%%%%%%%%%%%%%%%%%%%%%%%%%%%%%%%%%%%%%%%%%%%%%%%


\newglossaryentry{springbatch}{
  name=Spring Batch,
  description={Framework de Spring orientado al procesamiento por lotes (\textit{batch processing}), utilizado para ejecutar tareas asíncronas de larga duración.}
}

\newglossaryentry{flyway}{
  name=Flyway,
  description={Herramienta de migración de bases de datos que garantiza el versionado y la coherencia de los esquemas en los diferentes entornos.}
}

\newglossaryentry{postgresql}{
  name=PostgreSQL,
  description={Sistema de gestión de bases de datos relacional y de código abierto empleado para almacenar información estructurada.}
}

\newglossaryentry{elasticsearch}{
  name=Elasticsearch,
  description={Motor de búsqueda y análisis distribuido basado en documentos JSON. En la BIC Platform forma parte del ciclo de vida de un proceso, ya que almacena e indexa las instancias y variables generadas durante su ejecución en Process Execution, permitiendo consultar y analizar el estado, métricas y resultados de los procesos en tiempo real.}
}

\newglossaryentry{angular}{
  name=Angular,
  description={Framework de desarrollo web basado en TypeScript para la construcción de aplicaciones cliente (\textit{frontend}) modulares y escalables.}
}

\newglossaryentry{docker}{
  name=Docker,
  description={Plataforma de contenedores que permite empaquetar aplicaciones y sus dependencias en un entorno aislado y reproducible.}
}

\newglossaryentry{jira}{
  name=Jira,
  description={Herramienta de gestión ágil de proyectos que permite organizar tareas, historias de usuario y seguimiento de incidencias.}
}

\newglossaryentry{maven}{
  name=Maven,
  description={Herramienta de automatización y gestión de proyectos Java utilizada para la construcción, empaquetado y ejecución de pruebas.}
}

\newglossaryentry{aspose}{
  name=Aspose Cells,
  description={Biblioteca de Java empleada para la generación y manipulación avanzada de archivos en formato Excel.}
}

\newglossaryentry{kanban}{
  name=Kanban,
  description={Metodología ágil de gestión visual del trabajo basada en tableros y tarjetas, que permite controlar el flujo de tareas en tiempo real.}
}